\expandafter\def\csname HistoricalNoteText20\endcsname{Von der bedeutenden Discordanz zwischen verschiedenen Beobachtern bei einem und demselben Orte gibt schon das im vorhergehenden Artikel Mitgetheilte einige Proben; einige andere mögen hier noch angeführt werden, wo die Unterschiede viel grösser sind, als mit irgend einiger Wahrscheinlichkeit auf Rechnung regelmässiger jährlicher Änderung gesetzt werden kann. Die Inclination in Valparaiso war 1829 nach \textsc{King} $-40^0\,11'$, 1835 nach \textsc{Fitz Roy} $-38^0\,3'$. Auf der Insel Mauritius war die Intensität im Jahre 1818 nach \textsc{Freycinet} $1{,}096$, im Jahr 1836 nach \textsc{Fitz Roy} $1{,}192$. Noch grösser ist der Unterschied bei Otaheite, wo die Intensität 1830 von \textsc{Erman} $=1{,}172$ gefunden ist, hingegen 1835 von \textsc{Fitz Roy} $=1{,}017$. Diese letztere Verschiedenheit an einem für künftige Verbesserung der Elemente höchst wichtigen Platze ist bedeutend grösser, als die grösste, die unter allen unsern 86 Vergleichungen berechneter Intensitäten mit beobachteten vorkömmt.}
\SourceLeaf{166}\GaussVTheoryEarlyFolio{162}
\Article{30}
Wenn man bei der Beurtheilung der Unterschiede zwischen Rechnung und Beobachtung, welche die vorstehende tabellarische Vergleichung ergibt, in Erwägung zieht, dass einerseits fast sämmtliche Beobachtungen mit den Fehlern der Operation und den zufälligen Anomalien in der magnetischen Kraft selbst behaftet sind, und nicht für ein und dasselbe Jahr gelten\HistoricalNote{20}; andererseits, dass in unsern Formeln nur die Glieder bis zur vierten Ordnung enthalten sind, während die folgenden noch sehr merklich sein mögen: so scheint die Übereinstimmung zwischen Rechnung und Beobachtung allen billigen Erwartungen zu genügen, die man von einem ersten Versuche haben durfte. Unser Ausdruck für $\frac VR$ darf also wohl als der Wahrheit nahe kommend betrachtet werden, wenigstens in seinen beträchtlichern Gliedern, und es hat daher der Mühe werth geschienen, von dem Gange der numerischen Werthe von $\frac VR$ durch eine graphische Darstellung eine Versinnlichung zu geben. Es ist diess durch eine von Hrn. Dr. \textsc{Goldschmidt} gezeichnete Karte in drei Abtheilungen geschehen, deren erste nach \textsc{Mercator}’s Projection den ganzen Erdgürtel zwischen $70^0$ nordlicher und $70^0$ südlicher Breite, die beiden andern nach stereographischer Projection die Polargegenden bis zu $65^0$ Breite vorstellen. Die Correctionen und Vervollständigungen, welche in Zukunft eine wiederholte und auf vollkommenere Data gegründete Berechnung an dem Ausdruck für $\frac VR$ nöthig machen wird, werden zwar ohne Zweifel noch bedeutende Verschiebungen in diesem Liniensystem hervorbringen, besonders in den hohen südlichen Breiten: aber eine wesentliche Änderung in der ganzen Gestaltung selbst ist nicht denkbar ohne so grosse Änderungen in dem Ausdrucke für $\frac VR$, dass die Übereinstimmung mit den vorhandenen Beobachtungen verloren gehen müsste. Wir sind also hiedurch zu dem wichtigen Resultate geführt, dass das System der
\par\clearpage
\SourceLeaf{167}\GaussVTheoryEarlyFolio{163}
\noindent Linien gleicher Werthe von $V$ auf der Oberfläche der Erde wirklich unter dem einfachsten oben Art. 11 beschriebenen Typus begriffen ist, und dass also \textit{nur zwei magnetische Pole} auf der Erde vorhanden sind, wenn man von dem im 13. Artikel erwähnten Falle einer localen Ausnahme absieht, dessen Vorkommen oder Nichtvorkommen zur Zeit noch dahin gestellt bleiben muss. Die genaue Berechnung nach unsern Elementen gibt die Plätze dieser beiden Pole

1) in $73^0\,35'$ nordlicher Breite, $264^0\,21'$ Länge östlich von Greenwich, mit dem Werthe der ganzen Intensität $=1{,}701$ (nach gewöhnlicher Einheit).

2) in $72^0\,35'$ südlicher Breite, $152^0\,30'$ Länge mit der ganzen Intensität $=2{,}253$.

Im erstern Punkte hat $\frac VR$ seinen grössten Werth $=+895{,}86$, im zweiten den kleinsten $=-1030{,}24$.

Nach \textsc{Ross}’s Beobachtung fällt der nordliche magnetische Pol um $3^0\,30'$ südlicher als nach unserer Rechnung, und letztere gibt, wie aus unsrer Vergleichungstafel ersichtlich ist, eine um $1^0\,12'$ fehlerhafte Richtung der magnetischen Kraft an jenem Platze. Beim südlichen magnetischen Pole wird man eine bedeutend grössere Verschiebung zu erwarten haben. Da in Hobarttown, als dem demselben am nächsten liegenden Beobachtungsorte, die berechnete Inclination ohne Rücksicht auf das Zeichen, von der Rechnung um $3^0\,38'$ zu klein angegeben wird, insofern man sich auf die Beobachtung verlassen kann, so wird der wirkliche südliche magnetische Pol wahrscheinlich bedeutend nordlicher liegen als ihn unsere Rechnung angibt, und möchte derselbe etwa in der Gegend von $66^0$ Breite und $146^0$ Länge zu suchen sein.

\Article{31}
Wenngleich man den beiden Punkten auf der Erdoberfläche, wo die horizontale Kraft verschwindet, und die man die magnetischen Pole nennt, wegen ihrer Beziehung auf die Gestaltung der Erscheinungen der horizontalen Kraft auf der ganzen Erdfläche eine gewisse Bedeutsamkeit wohl beilegen mag, so muss man sich doch hüten, dieser Bedeutsamkeit eine weitere Ausdehnung zu geben: namentlich ist die Chorde, welche jene beiden Punkte verbindet, ohne alle Bedeutung, und es würde ein unpassender Missgriff sein, wenn man \textit{diese} gerade Linie durch die Benennung \textit{magnetische Axe} der Erde auszeichnen wollte. Die einzige Art, wie man dem Begriffe der magnetischen Axe eines Körpers eine all-
\GaussVTheoryGathering{21\textsuperscript{*}}
\par\clearpage
\SourceLeaf{168}\GaussVTheoryEarlyFolio{164}
\noindent gemein gültige Haltung geben kann, ist die im 5. Artikel der \textit{Intensitas vis magneticae} festgesetzte, wonach darunter eine gerade Linie verstanden wird, in Beziehung auf welche das Moment des in dem Körper enthaltenen freien Magnetismus ein Maximum ist. Zur Bestimmung der Lage der magnetischen Axe der Erde in diesem Sinn, und zugleich des Moments des Erdmagnetismus in Beziehung auf dieselbe, ist nun nach dem, was oben im 17. Art. bereits bemerkt ist, bloss die Kenntniss der Glieder erster Ordnung von $V$ erforderlich. Nach unsern Elementen Art. 26 ist
\[
P'=+925{,}782\cos u+89{,}024\sin u\cos\lambda-178{,}744\sin u\sin\lambda
\]
mithin sind $-925{,}782R^3$, $-89{,}024R^3$, $+178{,}744R^3$ die Momente des Erdmagnetismus in Beziehung auf die Erdaxe, und die beiden Erdradien für die Länge $0$ und $90^0$. Bei der Erdaxe ist die Richtung nach dem Nordpole zu verstanden, und das negative Zeichen des entsprechenden Moments zeigt an, dass die magnetische Axe einen stumpfen Winkel mit jener macht, d. i. dass ihr magnetischer Nordpol nach Süden gekehrt ist. Die Richtung der magnetischen Axe findet sich hieraus parallel dem Erddiameter von $77^0\,50'$ N. Breite $296^0\,29'$ Länge nach $77^0\,50'$ S. Breite $116^0\,29'$ Länge, und das magnetische Moment in Beziehung auf dieselbe $=947{,}08R^3$. Bei letzterm muss man sich erinnern, dass unsern Elementen eine Einheit für die Intensität zum Grunde liegt, die ein Tausendtheil der gewöhnlich gebrauchten ist. Um die Reduction auf die in der \textit{Intensitas vis magneticae} festgesetzte absolute Einheit zu erhalten, bemerken wir, dass in letzterer die horizontale Intensität in Göttingen, 1834 am 19. Julius $=1{,}7748$ gefunden war, woraus mit der Inclination $68^0\,1'$ die ganze Intensität $=4{,}7414$ folgt, während sie nach obiger Einheit $=1357$ angenommen wird. Der Reductionsfactor ist also $=0{,}0034941$, und sonach das magnetische Moment der Erde nach der absoluten Einheit
\[
=3{,}3092R^3
\]
Da bei dieser absoluten Einheit für die erdmagnetische Kraft das Millimeter als Längeneinheit angenommen ist, so muss auch $R$ in Millimetern angesetzt werden, wobei es, da ohnehin die Ellipticität der Erde hier nicht berücksichtigt wird, hinreichend ist, $R$ als Radius eines Kreises zu betrachten, dessen Umfang $40000$ Millionen Millimeter beträgt. Hienach wird obiges magnetische Moment durch
\expandafter\def\csname HistoricalNoteText21\endcsname{Insofern wir nemlich nicht befugt sind, bei allen magnetisirten Theilen der Erde durchaus parallele magnetische Axen vorauszusetzen. Je mehr an solchem Parallelismus fehlt, desto stärker muss die durchschnittliche Magnetisirung der Theile sein, um dasselbe magnetische Totalmoment hervorzubringen.}
\par\clearpage
\SourceLeaf{169}\GaussVTheoryEarlyFolio{165}
\noindent eine Zahl ausgedrückt, deren Logarithme $=29{,}93136$ oder durch $853800$ Quadrillionen. Nach derselben absoluten Einheit wurde das magnetische Moment eines einpfündigen Magnetstabes nach den im Jahre 1832 angestellten Versuchen $=100877000$ gefunden (\textit{Intensitas} Art. 21); das magnetische Moment der Erde ist also $8464$ Trillionen mal grösser. Es wären daher $8464$ Trillionen solcher Magnetstäbe, mit parallelen magnetischen Axen, erforderlich, um die magnetische Wirkung der Erde im äussern Raume zu ersetzen, was bei einer gleichförmigen Vertheilung durch den ganzen körperlichen Raum der Erde beinahe acht Stäbe (genauer $7{,}831$) auf jedes Kubikmeter beträgt. So ausgesprochen, behält dies Resultat seine Bedeutung, auch wenn man die Erde nicht als einen wirklichen Magnet betrachten, sondern den Erdmagnetismus blossen beharrlichen galvanischen Strömen in der Erde zuschreiben wollte. Betrachten wir aber die Erde als einen wirklichen Magnet, so sind wir genöthigt, \textit{durchschnittlich} wenigstens\HistoricalNote{21} jedem Theile derselben, der ein Achtel Kubikmeter gross ist, eine eben so starke Magnetisirung beizulegen, als jener Magnetstab enthält, ein Resultat, welches wohl den Physikern unerwartet sein wird.

\Article{32}
Die Art der wirklichen Vertheilung der magnetischen Flüssigkeiten in der Erde bleibt nothwendigerweise unbestimmt. In der That kann nach einem allgemeinen Theorem, welches bereits in der \textit{Intensitas} Art. 2 erwähnt ist, und bei einer andern Gelegenheit ausführlich behandelt werden soll, anstatt jeder beliebigen Vertheilung der magnetischen Flüssigkeiten innerhalb eines körperlichen Raumes allemal substituirt werden eine Vertheilung auf der Oberfläche dieses Raumes, so dass die Wirkung in jedem Punkte des äussern Raumes genau dieselbe bleibt, woraus man leicht schliesst, dass \textit{einerlei} Wirkung im ganzen äussern Raume aus unendlich vielen \textit{verschiedenen} Vertheilungen der magnetischen Flüssigkeiten im Innern abzuleiten ist.

Dagegen können wir diejenige fingirte Vertheilung auf der Oberfläche der Erde, welche der wirklichen im Innern, in Beziehung auf die daraus nach Aussen entstehenden Kräfte, vollkommen äquivalirt, angeben, und sogar, wegen der Ku-
\par\clearpage
\SourceLeaf{170}\GaussVTheoryEarlyFolio{166}
\noindent gelgestalt der Erde, auf eine höchst einfache Art. Es wird nemlich die Dichtigkeit des magnetischen Fluidums in jedem Punkte der Erdoberfläche. d. i. das Quantum des Fluidums, welches der Flächeneinheit entspricht, durch die Formel
\[
\frac1{4\pi}(\frac VR-2Z)
\]
ausgedrückt, oder durch
\[
-\frac1{4\pi}(3P'+5P''+7P'''+9P''''+\text{u. s. w.})
\]
Der Werth dieser Formel wird demnächst durch eine graphische Darstellung versinnlicht werden; hier mag nur bemerkt werden, dass er negativ an der nordlichen, positiv an der südlichen Hälfte der Erde ist, so jedoch, dass die Scheidungslinie den Äquator zweimal schneidet (in $6^0$ und $186^0$ Länge) und sich auf beiden Seiten bis zu etwa $15^0$ nordlicher und südlicher Breite von demselben entfernt; ferner dass auf der nordlichen Hälfte \textit{zwei} Minima Statt finden, auf der südlichen hingegen nur ein Maximum. Nach einer flüchtigen Rechnung finden sich diese Minima und das Maximum
\begin{center}
\begin{tabular}{rlrlrl}
$-209{,}1$&in&$55^0$&N. Breite&$263^0$&Länge\\
$-200{,}0$&in&$71^0$&N. Breite&$116^0$&Länge\\
$+277{,}7$&in&$70^0$&S. Breite&$154^0$&Länge
\end{tabular}
\end{center}
Bei den Werthen selbst liegt die Einheit unsrer Elemente zum Grunde, und sie müssen daher noch mit $0{,}0034941$ multiplicirt werden, wenn sie in absolutem Maass ausgedrückt werden sollen.

\Article{33}
Unsere Elemente sollen, wie schon oben bevorwortet ist, für nichts weiter gelten, als für eine erste Annäherung, und als solche stimmen sie nach Art. 29 mit den Beobachtungen befriedigend genug überein. Es leidet keinen Zweifel, dass eine Verbesserungsrechnung nach diesen Beobachtungen eine viel grössere Übereinstimmung verschaffen würde, und eine solche Rechnung würde an sich weiter keine Schwierigkeit haben als ihre Länge, die immer noch abschreckend gross bleibt, auch wenn man zur Abkürzung ähnliche Kunstgriffe anwenden wollte, wie von den Astronomen bei Verbesserung der Elemente der Planeten- und Kometenbahnen benutzt werden. Obgleich indessen diese Schwierigkeit leicht
\par\clearpage
\SourceLeaf{171}\GaussVTheoryEarlyFolio{167}
\noindent überwindlich sein würde, wenn die Arbeit unter eine Anzahl von Rechnern vertheilt werden könnte, so möchte es doch nicht gerathen sein, eine solche Verbesserung schon jetzt vorzunehmen, wo die Data von so vielen Plätzen, deren Mitbenutzung wesentlich sein würde, noch so geringe Zuverlässigkeit haben. Es wird am besten sein, vorerst die Vergleichung der Elemente mit Beobachtungen weiter fortzusetzen, wodurch man das Mittel finden wird, den allgemeinen Karten eine viel grössere Zuverlässigkeit zu geben, als bei dem bisher ausschliesslich empirischen Verfahren möglich war. Es sei uns aber erlaubt, einige Blicke auf die künftigen Fortschritte der Theorie zu werfen, deren völlige Realisirung freilich noch sehr entfernt sein mag.

\Article{34}
Zu einer befriedigenden Ausfeilung und Vervollständigung der Elemente müssen an die Beobachtungsdata viel höhere Forderungen gemacht werden, als bisher erfüllt sind. Jene sollten an allen zu benutzenden Punkten eine Schärfe haben, die bisjetzt nur an äusserst wenigen erreicht ist; sie sollten von den unregelmässigen Bewegungen gereinigt sein; sie sollten für Einerlei Zeitpunkt gelten. Es wird noch lange dauern, bis solchen Forderungen genügt werden kann: was aber zunächst am meisten Noth thut, ist die Herbeischaffung von \textit{vollständigen} (d. i. alle drei Elemente umfassenden) Beobachtungen an einem oder dem andern Punkte innerhalb derjenigen grossen Flächenräume, wo dergleichen bisher noch ganz fehlen; denn in der That hat ein neu hinzukommender Punkt allemal für die allgemeine Theorie desto grössere Wichtigkeit, je weiter er von den andern schon zu unserm Besitz gehörenden entfernt liegt.

Nach einer hinlänglichen Zwischenzeit wird man für einen zweiten Zeitpunkt die Elemente von neuem bestimmen, und so ihre Säcularänderungen erhalten. Aber offenbar wird dazu unumgänglich nöthig sein, das bisherige Maass der Intensitäten ganz fahren zu lassen, und ein absolutes an dessen Stelle zu setzen.

Im Laufe künftiger Jahrhunderte werden auch diese Änderungen nicht mehr als gleichförmig erscheinen, und die Erforschung des Ganges, in dem die Elemente fortschreiten, wird den Naturforschern unerschöpflichen Stoff zu Untersuchungen darbieten.
\par\clearpage
